% Generated draft: requires mathematical and source review.
\documentclass[11pt]{article}
\usepackage[margin=25mm]{geometry}
\usepackage{fontspec}
\setmainfont{Latin Modern Roman}
\usepackage{amsmath,amssymb,mathtools,unicode-math}
\setmathfont{Latin Modern Math}
\usepackage{enumitem,needspace,xurl,hyperref}
\setlength{\parindent}{0pt}
\setlength{\parskip}{5pt}
\setlength{\emergencystretch}{3em}
\newcommand{\N}{\mathbb{N}}
\newcommand{\Z}{\mathbb{Z}}
\newcommand{\Q}{\mathbb{Q}}
\newcommand{\R}{\mathbb{R}}
\newcommand{\C}{\mathbb{C}}
\newcommand{\F}{\mathbb{F}}
\newcommand{\A}{\mathbb{A}}
\newcommand{\PP}{\mathbb{P}}
\newcommand{\E}{\mathbb{E}}
\newcommand{\eps}{\varepsilon}
\newcommand{\dd}{\,\mathrm{d}}
\begin{document}
\section*{UnsolvedMath ID: 30006655}
\textbf{Source problem code:} catalogue entry
\subsection{Definitions and mathematical statement}
Let \(G\) be a finite group and let \(p\) be a prime. Under the usual convention, put
\[
\mathcal S_p(G)=\{P\leq G: P\neq 1,\ |P|=p^a\text{ for some integer }a\geq 1\},
\]
ordered by inclusion. Its order complex \(\Delta(\mathcal S_p(G))\) has as simplices the nonempty strict chains of these subgroups, with \(G\) acting by conjugation. Write \(K_p(G)=|\Delta(\mathcal S_p(G))|\) for its geometric realization. The supplied excerpt does not define \(\mathcal S_p(G)\); the nonidentity-subgroup convention must therefore be confirmed against the source. In particular, if \(p\nmid |G|\), this convention gives the empty complex, and the claims below retain this boundary case.

For \(G\)-spaces \(X,Y\), write \(X\simeq_G Y\) when there are continuous \(G\)-equivariant maps in both directions whose composites are homotopic to the respective identity maps through \(G\)-equivariant homotopies.

The supplied record does not define a synthetic building, minimality, a nontrivial synthetic building, or a homomorphic image of the subgroup complex. Accordingly, let \(\mathscr B_{\min}(G,p)\) be a placeholder for the source's intended collection of minimal synthetic buildings for \((G,p)\), regarded as \(G\)-spaces. Let \(\mathscr I(G,p)\) be a placeholder for its intended collection of homomorphic images of \(\Delta(\mathcal S_p(G))\), with the intended induced \(G\)-action. These are placeholders for missing definitions, not independent definitions of those objects. The precise category and meaning of image must be supplied before the image assertion is fully determined.

Subject to those unresolved definitions, the supplied problem asks for the following assertions.
\begin{enumerate}
\item For every finite group \(G\) and every prime \(p\), there are only finitely many \(G\)-equivariant homotopy types represented by \(\mathscr B_{\min}(G,p)\):
\[
\#\bigl(\mathscr B_{\min}(G,p)/{\simeq_G}\bigr)<\infty.
\]
This finiteness assertion occurs in the supplied statement and clean statement, but not in the supplied original excerpt containing Conjectures 5--7 and Theorem 8. Its precise source location and any additional hypotheses remain unprovided.
\item \textbf{Conjecture 5.} For every finite group \(G\) and every prime \(p\),
\[
\exists X\in\mathscr B_{\min}(G,p)\quad K_p(G)\simeq_G X.
\]
\item \textbf{Conjecture 6.} For every finite group \(G\), every prime \(p\), and every \(X\in\mathscr B_{\min}(G,p)\),
\[
\exists Y\in\mathscr I(G,p)\quad X\simeq_G Y.
\]
The term ``homomorphic image'' is retained as an unresolved source term; no choice of quotient, simplicial image, or other image construction is made here.
\item \textbf{Conjecture 7.} If \(G\) is a finite group of Lie type in characteristic \(p\), then every \(X\in\mathscr B_{\min}(G,p)\) that is nontrivial in the source's sense satisfies
\[
X\simeq_G T(G),
\]
where \(T(G)\) is the Tits building referred to in the source, with its natural \(G\)-action. The excerpt does not specify the precise convention for groups of Lie type, their building, or the meaning of nontrivial; in particular, nontrivial is not silently identified with noncontractible.
\end{enumerate}

\textbf{Reported partial result and status qualification.} The supplied original excerpt states as Theorem 8 that Conjecture 7 holds for groups of Lie type of Lie rank at most \(2\). The supplied literature assessment dated 2026-08-22 describes the problem as partially solved: it reports that rank-at-most-two result and describes Conjectures 5 and 6 and the general Lie-rank classification as unresolved. These are supplied source and curation claims, not independently verified results or a new assessment. The assessment does not separately resolve the additional finiteness assertion.
\subsection{Short English statement}
For each finite group \(G\) and prime \(p\), the supplied problem asks whether minimal synthetic buildings have finitely many \(G\)-equivariant homotopy types, whether the \(p\)-subgroup complex represents one of them, and whether every such building is represented by a homomorphic image of that complex. For groups of Lie type in characteristic \(p\), it further asks whether every nontrivial example has the equivariant homotopy type of the Tits building. Essential source definitions are missing, and the finiteness clause is absent from the supplied original excerpt. The source reports the Lie-rank-at-most-two case as proved; the supplied dated status remains subject to independent review.
\subsection{Sources}
\begin{itemize}
\item Local registry: \texttt{data/unsolvedmath/problems.json}, supplied record with permanent UnsolvedMath ID \(30006655\), title ``Classification of Minimal Synthetic Buildings,'' source problem code \texttt{catalogue entry}. The record gives proposed year 2026 and no named proposer.
\item \emph{Cohomology of Finite Groups: Interactions and Applications} (2026), Oberwolfach Reports, DOI \texttt{10.4171/owr/2026/9}, pp.\ 541--610. The supplied original excerpt identifies p.\ 583 and contains Conjectures 5--7 and Theorem 8. Supplied source URL and literature-source link: \url{https://doi.org/10.4171/owr/2026/9}.
\item Gullerud--Webb, \emph{Synthetic Buildings for Finite Groups}. Supplied literature-source link: \url{https://arxiv.org/abs/2608.13414}. The record provides this reference without a theorem number or a passage establishing its relationship to each assertion.
\item The exact supplied literature-check date is \texttt{2026-08-22}. Its assessment reports a partially solved status, the rank-at-most-two case of Conjecture 7 as proved, and Conjectures 5 and 6 and the general Lie-rank classification as unresolved. This is supplied dated curation derived from the Oberwolfach record, not independently verified literature assessment.
\item The registry labels the statement \texttt{corrected\_verified} and says it was reconstructed and checked against the official Oberwolfach report. That verification claim is supplied metadata; no source was independently visited or checked in preparing this draft.
\end{itemize}
\end{document}
